%! AP Statistics — Unit 7, Lesson 7.3 — Homework (Answer Key)
\documentclass[10pt]{article}
\usepackage{apstats-article}
\usepackage{apstats-key}
\newcommand{\TallMath}[1]{$\displaystyle #1\rule[-1.4em]{0pt}{3.2em}$}

\begin{document}

\pageheader{Unit 7, Lesson 7.3}{Homework --- Answer Key}
\namedateperiod

\begin{enumerate}

  \item \textbf{Interpretation Practice.} 90\% CI: $(3.8, 5.6)$ hr, $n = 45$.

        \begin{enumerate}[label=\alph*.]
          \item Correct interpretation:
                \ansline{We are 90\% confident that the interval (3.8, 5.6) hours captures the true mean number of hours per week adults exercise.}

          \item Claim: mean $= 6$ hr. Evidence against?
                \ansline{Yes: 6 hours lies above the entire interval (3.8, 5.6). The data provide convincing evidence that the true mean exercise time differs from (specifically, is less than) 6 hours per week.}

          \item Claim: mean $< 5$ hr. Supported?
                \ansline{Not convincingly: the interval (3.8, 5.6) includes values both below and above 5 hours. Because the upper bound exceeds 5, we cannot conclude that the mean is less than 5. The data do not provide convincing evidence that the mean is below 5 hours.}
        \end{enumerate}

  \item \textbf{Width Relationships.} $\bar x = \$87.40$, $s = \$24.00$, $n = 36$, 95\% CI.

        \begin{enumerate}[label=\alph*.]
          \item $df = $ \ans{35}, $t^* = $ \ans{2.030}, $SE = 24/\sqrt{36} = $ \ans{\$4.00}

                $ME = 2.030 \times 4.00 = $ \ans{\$8.12}

                Interval: $87.40 \pm 8.12 = ($ \ans{\$79.28}$,$ \ans{\$95.52}$)$

          \item For $n = 144$: $SE = 24/\sqrt{144} = $ \ans{\$2.00}.
                \ansline{$ME = 2.030 \times 2.00 = \$4.06$. The margin of error is halved (from \$8.12 to \$4.06), because quadrupling $n$ halves $1/\sqrt{n}$.}

          \item 99\% CI ($n = 36$): $t^*$ increases (from 2.030 to $\approx 2.724$), so the interval is \ans{wider}. \ansline{$ME = 2.724 \times 4.00 = \$10.90$; interval $\approx (\$76.50, \$98.30)$ --- wider than the 95\% CI.}
        \end{enumerate}

  \item \textbf{AP-Style.} $n = 22$, $\bar x = 14.6$ mi, $s = 8.3$ mi, 95\% CI.

        \begin{enumerate}[label=\alph*.]
          \item Full procedure:
                \ansline{Parameter: $\mu = $ true mean commute distance (miles) for all employees at the company.}
                \ansline{Conditions: Random: random sample of employees. \checkmark \quad 10\%: $22 \le 10\%$ of all employees (assume $N \ge 220$). \checkmark \quad Large Sample: $n = 22 < 30$; must check for no strong skewness or outliers in the data. Assume verified. \checkmark}
                \ansline{Calculate: $df = 21$; $t^* = 2.080$; $SE = 8.3/\sqrt{22} = 1.770$ mi. $ME = 2.080 \times 1.770 = 3.682$ mi.}

                Interval: $14.6 \pm 3.682 = ($ \ans{10.918}$,$ \ans{18.282}$)$ miles

                \ansline{Conclude: We are 95\% confident that the interval (10.9, 18.3) miles captures the true mean commute distance for all employees at this company.}

                \begin{teachernote}
                    $t^*$ for $df = 21$, 95\% CI is 2.080. $SE = 8.3/\sqrt{22} \approx 1.770$.
                \end{teachernote}

          \item Claim: mean $> 12$ mi?
                \ansline{The interval (10.9, 18.3) miles includes values below 12 miles (the lower bound is $\approx 10.9$), so the data do not provide convincing evidence that the true mean commute exceeds 12 miles. The claimed value of 12 miles is inside the interval.}
        \end{enumerate}

\end{enumerate}

\begin{reflectionbox}
\textbf{Reflection:}
\ansline{Disagree: a 99\% CI is more reliable (captures $\mu$ more often across repeated samples), but it is also wider. The trade-off is precision: higher confidence levels produce larger margins of error, giving a less informative range of plausible values. A 95\% CI is narrower and thus more useful for decision-making, at the cost of being wrong 5\% of the time rather than 1\%.}
\end{reflectionbox}

\end{document}
